\chapter{Bulk Reduced-Model Classes}
\index{model class!bulk|textbf}

\label{ch:model-reduction}

\chapterstatus{proposed bulk-reduction framework; no fitted or selected
bulk tight-binding model is claimed.}

Reduction replaces an accepted parent representation by a member of a more
restricted model class for a declared purpose. It is not merely a change of
basis.

\section{Bulk reduction program}

\subsection{Reduction as a constrained problem}
\index{reduction!constrained problem}


Let $H_{\mathrm p}$ be a compatible finite representation of a parent operator
and let
\begin{equation}
  \mathfrak M_j
  =
  \{H_j(\boldsymbol\theta):\boldsymbol\theta\in\Theta_j\}
\end{equation}
denote a model class indexed by $j$. The parameter domain $\Theta_j$ includes
all symmetry, range, orbital, and sign conventions imposed by that class. A
reduction operation must state

\begin{itemize}
  \item the parent operator and its retained representation;
  \item the candidate model class and parameter constraints;
  \item the fitting or projection map;
  \item the training data and objective;
  \item the withheld validation data;
  \item the metrics and tolerances associated with the intended use; and
  \item the provenance of the resulting parameters and operator.
\end{itemize}

The least complex model is meaningful only relative to a nested or otherwise
ordered family and a complete acceptance rule. Fewer parameters do not make a
model acceptable when one required criterion fails.

\subsection{Bulk orthogonal \texorpdfstring{$sp^3s^{\ast}$}{sp3s*} hierarchy}
\index{tight-binding model!orthogonal sp3s-star}
\index{model hierarchy!bulk}


The initial bulk lattice model is an orthogonal Slater--Koster
$sp^3s^{\ast}$ model for diamond-structure
silicon~\cite{slater1954,vogl1983}. Each silicon site carries a
prescribed orbital set, and symmetry relates hopping matrix elements to a
smaller parameter vector. The executable operator must define site, orbital,
spin, neighbor-shell, lattice, Fourier, and energy-zero conventions before any
parameter value is interpreted.

The first model class uses the approved nearest-neighbor structure. If that
class cannot satisfy the declared criteria, later classes may introduce
second-neighbor or selected symmetry-allowed corrections in a fixed sequence.
The sequence must be frozen before the compatibility search so that model
selection is not changed after inspecting the outcomes.

Hermiticity, crystal symmetries, and controlled limiting cases are
software-verification requirements for every model class. Passing them shows
that the implementation follows its mathematical contract; it does not show
that the class is expressive enough for the silicon parent.

\subsection{Two bulk reduction routes}
\index{reduction route!direct}
\index{reduction route!Wannier-mediated}


Stage~04 compares two paths that share the same accepted Kohn--Sham parent:
\begin{align}
  \text{direct path:}\quad
  &\mathrm{DFT}\longrightarrow\mathrm{TB},\\
  \text{operator-mediated path:}\quad
  &\mathrm{DFT}\longrightarrow\mathrm{Wannier}
    \longrightarrow\mathrm{TB}.
\end{align}
The direct path fits selected spectral and band-edge information. The mediated
path first constructs a validated Wannier operator and then aligns the
Wannier and tight-binding orbital spaces so that represented operator residuals
can be defined.

The paths cannot be compared merely because both produced parameter vectors.
Their result records must identify the same accepted parent dataset and source
manifest and must carry compatible physical, numerical, pseudopotential,
workflow, energy, representation, artifact, and validation metadata. This
provenance-compatible join prevents path differences from being attributed to
reduction when the branches actually consumed different parents.

\subsection{Training and withheld validation}
\index{training data!bulk model}
\index{withheld validation}


Training data determine or constrain $\boldsymbol\theta$; withheld data test the
resulting model after fitting. The two sets must be frozen and disjoint. For the
bulk pilot, validation includes withheld band energies, the indirect gap,
conduction-valley position, and longitudinal and transverse electron effective
masses. Operator metrics are added only where a common aligned representation
exists.

A weighted spectral objective may be written abstractly as
\begin{equation}
  \mathcal L_E(\boldsymbol\theta)
  =
  \sum_{q\in\mathcal T_E}
  w_q\,
  \ell_q\!\left(
    H_j(\boldsymbol\theta),H_{\mathrm p}
  \right),
\end{equation}
where $\mathcal T_E$ is the frozen spectral training set, $w_q$ are declared
weights, and every loss term $\ell_q$ defines its units and normalization. The
objective value on training data is not a withheld validation result.

\subsection{Aligned operator objective}
\index{objective!aligned operator}


Let $H_{\mathrm W}(\mathbf R)$ be the retained real-space Wannier blocks and
$H_{\mathrm{TB}}(\mathbf R;\boldsymbol\theta)$ the corresponding tight-binding
blocks. For a declared set $\mathcal U$ of symmetry-compatible alignment maps,
one possible operator objective is
\begin{equation}
  \mathcal L_H(C,\boldsymbol\theta)
  =
  \sum_{\mathbf R}\omega_{\mathbf R}
  \left\|
    H_{\mathrm W}(\mathbf R)
    -C H_{\mathrm{TB}}(\mathbf R;\boldsymbol\theta)C^{\dagger}
  \right\|_{\mathrm F}^{2},
  \qquad C\in\mathcal U.
\end{equation}
The direction of $C$, the retained $\mathbf R$ blocks, shell weights
$\omega_{\mathbf R}$, normalization, units, and treatment of omitted blocks
must be fixed. Optimizing over $C$ does not waive the subspace, geometry,
orbital, spin, or energy-reference compatibility requirements established in
Chapter~\ref{ch:representations-and-alignment}.

The residual can be decomposed by onsite and hopping contribution, orbital
block, symmetry channel, and neighbor shell. Such a decomposition may identify
where a model class lacks expressiveness, but it remains a reduction diagnostic
relative to the selected parent representation.

\subsection{Spectral and operator admissibility}
\index{admissibility!spectral}
\index{admissibility!operator}


For a model class $\mathfrak M_j$, define the spectral-admissible set
\begin{equation}
  \mathfrak A_{E,j}(\tau_E)
  =
  \left\{
    H_j(\boldsymbol\theta)\in\mathfrak M_j:
    \mathcal L_E(\boldsymbol\theta)\leq\tau_E
  \right\}
\end{equation}
and the operator-admissible set
\begin{equation}
  \mathfrak A_{H,j}(\tau_H)
  =
  \left\{
    H_j(\boldsymbol\theta)\in\mathfrak M_j:
    \min_{C\in\mathcal U}\mathcal L_H(C,\boldsymbol\theta)
    \leq\tau_H
  \right\}.
\end{equation}
The tolerances $\tau_E$ and $\tau_H$ are proposed acceptance inputs, not values
selected by this chapter. They must be fixed before the compatibility result is
interpreted.

A common acceptable Hamiltonian exists when
\begin{equation}
  \mathfrak A_{E,j}(\tau_E)
  \cap
  \mathfrak A_{H,j}(\tau_H)
  \neq\varnothing.
\end{equation}
If the intersection is empty, a declared dimensionless metric
$d_{\mathrm{TB}}$ on common canonical model coordinates may define the set
separation
\begin{equation}
  \delta_j^{\ast}
  =
  \inf_{\substack{
    H_E\in\mathfrak A_{E,j}\\
    H_H\in\mathfrak A_{H,j}}}
  d_{\mathrm{TB}}(H_E,H_H).
\end{equation}
A numerical optimizer returning two parameter vectors does not by itself prove
that this infimum has been found. Optimization stability, parameter
identifiability, class sensitivity, and the definition of the metric remain
part of the evidence.

The bulk model-selection rule is the lowest-order tested class satisfying all
prespecified criteria. If no tested class passes, the correct outcome is that
no accepted reduced class was found in the tested hierarchy; complexity is not
expanded after the fact merely to force a positive result.

\section{Path dependence and error accounting}

\subsection{Path dependence and noncommutativity}
\index{path dependence}
\index{noncommutativity!reduction}


Reduction and impurity extraction need not commute. The program therefore
compares
\begin{align}
  &\text{extract then reduce},\\
  &\text{reduce parent operators then extract}.
\end{align}
After aligning their outputs, a declared norm measures the defect between the
two paths. Similar defects compare direct and Wannier-mediated bulk reductions
and test gauge equivariance.

A nonzero defect is not automatically an implementation bug. It may arise from
projection, fitting, truncation, incompatible approximation domains, or
sensitivity to alignment. The final interpretation requires precisely defined
paths, compatible outputs, a stated norm, a validation domain, and evidence
that separates numerical artifacts from structural noncommutativity.

\subsection{Error accounting}
\index{error!accounting}


Every reduction result remains relative to the chosen Kohn--Sham parent.
Parent-model error concerns that parent's relationship to independent physical
references. Numerical error concerns approximation of the parent calculation
and the reduction algorithms. Model-reduction error concerns replacement by a
restricted class. Impurity finite-size error and continuum discretization may
require still more specific components.

No scalar ``total error'' is reported until an owning contract defines how its
components relate and whether their units, norms, correlations, and validation
domains are compatible. The immediate purpose of the hierarchy is instead to
make each loss visible and to determine which scientific statements survive a
particular sequence of reductions.
